% Options for packages loaded elsewhere
% Options for packages loaded elsewhere
\PassOptionsToPackage{unicode}{hyperref}
\PassOptionsToPackage{hyphens}{url}
\PassOptionsToPackage{dvipsnames,svgnames,x11names}{xcolor}
%
\documentclass[
  authoryear,
  preprint]{elsarticle}
\usepackage{xcolor}
\usepackage{amsmath,amssymb}
\setcounter{secnumdepth}{5}
\usepackage{iftex}
\ifPDFTeX
  \usepackage[T1]{fontenc}
  \usepackage[utf8]{inputenc}
  \usepackage{textcomp} % provide euro and other symbols
\else % if luatex or xetex
  \usepackage{unicode-math} % this also loads fontspec
  \defaultfontfeatures{Scale=MatchLowercase}
  \defaultfontfeatures[\rmfamily]{Ligatures=TeX,Scale=1}
\fi
\usepackage{lmodern}
\ifPDFTeX\else
  % xetex/luatex font selection
\fi
% Use upquote if available, for straight quotes in verbatim environments
\IfFileExists{upquote.sty}{\usepackage{upquote}}{}
\IfFileExists{microtype.sty}{% use microtype if available
  \usepackage[]{microtype}
  \UseMicrotypeSet[protrusion]{basicmath} % disable protrusion for tt fonts
}{}
\makeatletter
\@ifundefined{KOMAClassName}{% if non-KOMA class
  \IfFileExists{parskip.sty}{%
    \usepackage{parskip}
  }{% else
    \setlength{\parindent}{0pt}
    \setlength{\parskip}{6pt plus 2pt minus 1pt}}
}{% if KOMA class
  \KOMAoptions{parskip=half}}
\makeatother
% Make \paragraph and \subparagraph free-standing
\makeatletter
\ifx\paragraph\undefined\else
  \let\oldparagraph\paragraph
  \renewcommand{\paragraph}{
    \@ifstar
      \xxxParagraphStar
      \xxxParagraphNoStar
  }
  \newcommand{\xxxParagraphStar}[1]{\oldparagraph*{#1}\mbox{}}
  \newcommand{\xxxParagraphNoStar}[1]{\oldparagraph{#1}\mbox{}}
\fi
\ifx\subparagraph\undefined\else
  \let\oldsubparagraph\subparagraph
  \renewcommand{\subparagraph}{
    \@ifstar
      \xxxSubParagraphStar
      \xxxSubParagraphNoStar
  }
  \newcommand{\xxxSubParagraphStar}[1]{\oldsubparagraph*{#1}\mbox{}}
  \newcommand{\xxxSubParagraphNoStar}[1]{\oldsubparagraph{#1}\mbox{}}
\fi
\makeatother


\usepackage{longtable,booktabs,array}
\usepackage{calc} % for calculating minipage widths
% Correct order of tables after \paragraph or \subparagraph
\usepackage{etoolbox}
\makeatletter
\patchcmd\longtable{\par}{\if@noskipsec\mbox{}\fi\par}{}{}
\makeatother
% Allow footnotes in longtable head/foot
\IfFileExists{footnotehyper.sty}{\usepackage{footnotehyper}}{\usepackage{footnote}}
\makesavenoteenv{longtable}
\usepackage{graphicx}
\makeatletter
\newsavebox\pandoc@box
\newcommand*\pandocbounded[1]{% scales image to fit in text height/width
  \sbox\pandoc@box{#1}%
  \Gscale@div\@tempa{\textheight}{\dimexpr\ht\pandoc@box+\dp\pandoc@box\relax}%
  \Gscale@div\@tempb{\linewidth}{\wd\pandoc@box}%
  \ifdim\@tempb\p@<\@tempa\p@\let\@tempa\@tempb\fi% select the smaller of both
  \ifdim\@tempa\p@<\p@\scalebox{\@tempa}{\usebox\pandoc@box}%
  \else\usebox{\pandoc@box}%
  \fi%
}
% Set default figure placement to htbp
\def\fps@figure{htbp}
\makeatother





\setlength{\emergencystretch}{3em} % prevent overfull lines

\providecommand{\tightlist}{%
  \setlength{\itemsep}{0pt}\setlength{\parskip}{0pt}}



 
\usepackage[]{natbib}
\bibliographystyle{elsarticle-harv}


\usepackage{fvextra}
\usepackage{witharrows}
\usepackage{amsmath}
\usepackage{listings}
\usepackage{geometry}[margin=0.25in]
\usetikzlibrary{calc}
\DefineVerbatimEnvironment{Highlighting}{Verbatim}{breaklines,commandchars=\\\{\}}
\lstset{
  basicstyle=\small\ttfamily,
  breaklines=true,
  frame=single,
  numbers=left,
  numberstyle=\tiny\color{gray},
  captionpos=b,
  language=python
}
\makeatletter
\@ifpackageloaded{caption}{}{\usepackage{caption}}
\AtBeginDocument{%
\ifdefined\contentsname
  \renewcommand*\contentsname{Table of contents}
\else
  \newcommand\contentsname{Table of contents}
\fi
\ifdefined\listfigurename
  \renewcommand*\listfigurename{List of Figures}
\else
  \newcommand\listfigurename{List of Figures}
\fi
\ifdefined\listtablename
  \renewcommand*\listtablename{List of Tables}
\else
  \newcommand\listtablename{List of Tables}
\fi
\ifdefined\figurename
  \renewcommand*\figurename{Figure}
\else
  \newcommand\figurename{Figure}
\fi
\ifdefined\tablename
  \renewcommand*\tablename{Table}
\else
  \newcommand\tablename{Table}
\fi
}
\@ifpackageloaded{float}{}{\usepackage{float}}
\floatstyle{ruled}
\@ifundefined{c@chapter}{\newfloat{codelisting}{h}{lop}}{\newfloat{codelisting}{h}{lop}[chapter]}
\floatname{codelisting}{Listing}
\newcommand*\listoflistings{\listof{codelisting}{List of Listings}}
\makeatother
\makeatletter
\makeatother
\makeatletter
\@ifpackageloaded{caption}{}{\usepackage{caption}}
\@ifpackageloaded{subcaption}{}{\usepackage{subcaption}}
\makeatother
\makeatletter
\definecolor{QuartoInternalColor4}{rgb}{0.00,0.45,0.15}
\definecolor{QuartoInternalColor8}{rgb}{0.00,0.40,0.79}
\definecolor{QuartoInternalColor5}{rgb}{0.00,0.40,0.00}
\definecolor{QuartoInternalColor9}{rgb}{0.70,0.49,0.07}
\definecolor{QuartoInternalColor2}{rgb}{0,0,0}
\definecolor{QuartoInternalColor6}{rgb}{0.20,0.40,0.40}
\definecolor{QuartoInternalColor10}{rgb}{0.60,0.00,1.00}
\definecolor{QuartoInternalColor3}{rgb}{0.15,0.56,0.56}
\definecolor{QuartoInternalColor1}{rgb}{0.70,0.17,0.19}
\definecolor{QuartoInternalColor7}{rgb}{0.74,0.74,0.74}
\makeatother
\journal{Laboratory of Computational Physics}
\usepackage{bookmark}
\IfFileExists{xurl.sty}{\usepackage{xurl}}{} % add URL line breaks if available
\urlstyle{same}
\hypersetup{
  pdftitle={Neutrino Oscillation Analysis with Machine Learning},
  pdfauthor={Álvaro Méndez Rodríguez de Tembleque; Miguel Orden; Adrián Orellana; Carlos Hita},
  colorlinks=true,
  linkcolor={blue},
  filecolor={Maroon},
  citecolor={Blue},
  urlcolor={Blue},
  pdfcreator={LaTeX via pandoc}}


\setlength{\parindent}{6pt}
\begin{document}

\begin{frontmatter}
\title{Neutrino Oscillation Analysis with Machine Learning}
\author[1]{Álvaro Méndez Rodríguez de Tembleque%
%
}
 \ead{alvaro.mendezrodriguezdetembleque@studenti.unipd.it} 
\author[1]{Miguel Orden%
%
}
 \ead{alvaro.mendezrodriguezdetembleque@studenti.unipd.it} 
\author[1]{Adrián Orellana%
%
}
 \ead{alvaro.mendezrodriguezdetembleque@studenti.unipd.it} 
\author[1]{Carlos Hita%
%
}
 \ead{alvaro.mendezrodriguezdetembleque@studenti.unipd.it} 

\affiliation[1]{organization={Universita degli Studi di
Padova},addressline={Department of physics and astronomy ``Galileo
Galilei''},city={City},postcode={12345},postcodesep={}}

\cortext[cor1]{Corresponding author}




        
\begin{abstract}
We present a machine learning-based approach to identify neutrino
interactions from a dataset dominated by radioactive background. The
analysis focuses on Inverse Beta Decay (IBD) event selection, using
engineered features derived from event timing, spatial correlation, and
energy deposition. A classification model is trained on simulated data
and calibrated to account for discrepancies between training and
real-world distributions. Selected events are used to reconstruct the
neutrino energy spectrum and extract oscillation parameters.
\end{abstract}





\end{frontmatter}
    

\section{Introduction}\label{introduction}

Neutrinos are weakly interacting particles that provide insight into
fundamental physics beyond the Standard Model. Their detection is
challenging due to extremely low interaction rates and significant
background contamination.

This work focuses on:

\begin{itemize}
\tightlist
\item
  Identifying neutrino events via Inverse Beta Decay
\item
  Separating signal from dominant radioactive background
\item
  Reconstructing the neutrino energy spectrum
\item
  Extracting oscillation signatures
\end{itemize}

\begin{center}\rule{0.5\linewidth}{0.5pt}\end{center}

\section{Theoretical Background}\label{theoretical-background}

\subsection{Neutrino Detection}\label{neutrino-detection}

Neutrinos are detected through the Inverse Beta Decay (IBD) process:

\[
\bar{\nu}_e + p \rightarrow n + e^+
\]

This interaction produces a characteristic \textbf{prompt-delayed}
signal:

\begin{itemize}
\tightlist
\item
  Prompt: positron energy deposition\\
\item
  Delayed: neutron capture after a short time delay
\end{itemize}

\subsection{Signal and Background}\label{signal-and-background}

Signal events are characterized by:

\begin{itemize}
\tightlist
\item
  Two events close in time
\item
  Spatial proximity
\item
  Correlated energy signatures
\end{itemize}

Background arises primarily from random coincidences of radioactive
decays, which can mimic the signal.

\subsection{Neutrino Oscillations}\label{neutrino-oscillations}

Neutrino oscillation probabilities depend on the ratio:

\[
\frac{L}{E}
\]

where:

\begin{itemize}
\tightlist
\item
  \(L\) = distance from source to detector\\
\item
  \(E\) = neutrino energy
\end{itemize}

Oscillations are observed through distortions in the measured energy
spectrum.

\begin{center}\rule{0.5\linewidth}{0.5pt}\end{center}

\section{Dataset and Preprocessing}\label{dataset-and-preprocessing}

\subsection{Dataset Description}\label{dataset-description}

The dataset consists of sequential events, each characterized by:

\begin{itemize}
\tightlist
\item
  Energy
\item
  Position
\item
  Time :contentReference{oaicite:1}
\end{itemize}

Simulated data provides labeled examples of:

\begin{itemize}
\tightlist
\item
  Signal (IBD events)
\item
  Background events
\end{itemize}

\subsection{Data Loading}\label{data-loading}

\begin{codelisting}

\caption{\label{lst-load-data}}

\centering{

\begin{Highlighting}
\textcolor{black}{ModuleNotFoundError: No module named 'pkg\_resources'}
\textcolor{black}{}\textcolor{QuartoInternalColor1}{---------------------------------------------------------------------------}\textcolor{QuartoInternalColor2}{}
\textcolor{QuartoInternalColor2}{}\textcolor{QuartoInternalColor1}{ModuleNotFoundError}\textcolor{QuartoInternalColor2}{                       Traceback (most recent call last)}
\textcolor{QuartoInternalColor2}{}\textcolor{QuartoInternalColor3}{Cell}\textcolor{QuartoInternalColor2}{}\textcolor{QuartoInternalColor3}{ }\textcolor{QuartoInternalColor2}{}\textcolor{QuartoInternalColor4}{In[1]}\textcolor{QuartoInternalColor2}{}\textcolor{QuartoInternalColor4}{, line 2}\textcolor{QuartoInternalColor2}{}
\textcolor{QuartoInternalColor2}{}\textcolor{QuartoInternalColor4}{      1}\textcolor{QuartoInternalColor2}{ }\textcolor{QuartoInternalColor5}{import}\textcolor{QuartoInternalColor2}{ pyarrow.parquet }\textcolor{QuartoInternalColor5}{as}\textcolor{QuartoInternalColor2}{ pq}
\textcolor{QuartoInternalColor2}{}\textcolor{QuartoInternalColor4}{----> }\textcolor{QuartoInternalColor2}{}\textcolor{QuartoInternalColor4}{2}\textcolor{QuartoInternalColor2}{ }\textcolor{QuartoInternalColor5}{from}\textcolor{QuartoInternalColor2}{ ydata\_profiling }\textcolor{QuartoInternalColor5}{import}\textcolor{QuartoInternalColor2}{ ProfileReport}
\textcolor{QuartoInternalColor2}{}\textcolor{QuartoInternalColor4}{      3}\textcolor{QuartoInternalColor2}{ }\textcolor{QuartoInternalColor5}{import}\textcolor{QuartoInternalColor2}{ pandas }\textcolor{QuartoInternalColor5}{as}\textcolor{QuartoInternalColor2}{ pd}
\textcolor{QuartoInternalColor2}{}\textcolor{QuartoInternalColor4}{      4}\textcolor{QuartoInternalColor2}{ }
\textcolor{QuartoInternalColor2}{}\textcolor{QuartoInternalColor4}{      5}\textcolor{QuartoInternalColor2}{ }\textcolor{QuartoInternalColor6}{# Read as a Table}\textcolor{QuartoInternalColor2}{}
\textcolor{QuartoInternalColor2}{}\textcolor{QuartoInternalColor3}{File }\textcolor{QuartoInternalColor2}{}\textcolor{QuartoInternalColor4}{\textasciitilde{}/Desktop/master/2sem/LCPB/projects/.venv/lib/python3.13/site-packages/ydata\_profiling/\_\_init\_\_.py:11}\textcolor{QuartoInternalColor2}{}
\textcolor{QuartoInternalColor2}{}\textcolor{QuartoInternalColor4}{      8}\textcolor{QuartoInternalColor2}{ }\textcolor{QuartoInternalColor5}{import}\textcolor{QuartoInternalColor2}{}\textcolor{QuartoInternalColor7}{ }\textcolor{QuartoInternalColor2}{}\textcolor{QuartoInternalColor2}{importlib}\textcolor{QuartoInternalColor2}{}\textcolor{QuartoInternalColor2}{.}\textcolor{QuartoInternalColor2}{}\textcolor{QuartoInternalColor2}{util}\textcolor{QuartoInternalColor2}{  }\textcolor{QuartoInternalColor6}{# isort:skip # noqa}\textcolor{QuartoInternalColor2}{}
\textcolor{QuartoInternalColor2}{}\textcolor{QuartoInternalColor4}{      9}\textcolor{QuartoInternalColor2}{ }\textcolor{QuartoInternalColor5}{from}\textcolor{QuartoInternalColor2}{}\textcolor{QuartoInternalColor7}{ }\textcolor{QuartoInternalColor2}{}\textcolor{QuartoInternalColor2}{warnings}\textcolor{QuartoInternalColor2}{}\textcolor{QuartoInternalColor7}{ }\textcolor{QuartoInternalColor2}{}\textcolor{QuartoInternalColor5}{import}\textcolor{QuartoInternalColor2}{ warn}
\textcolor{QuartoInternalColor2}{}\textcolor{QuartoInternalColor4}{---> }\textcolor{QuartoInternalColor2}{}\textcolor{QuartoInternalColor4}{11}\textcolor{QuartoInternalColor2}{ }\textcolor{QuartoInternalColor5}{from}\textcolor{QuartoInternalColor2}{}\textcolor{QuartoInternalColor7}{ }\textcolor{QuartoInternalColor2}{}\textcolor{QuartoInternalColor2}{ydata\_profiling}\textcolor{QuartoInternalColor2}{}\textcolor{QuartoInternalColor2}{.}\textcolor{QuartoInternalColor2}{}\textcolor{QuartoInternalColor2}{compare\_reports}\textcolor{QuartoInternalColor2}{}\textcolor{QuartoInternalColor7}{ }\textcolor{QuartoInternalColor2}{}\textcolor{QuartoInternalColor5}{import}\textcolor{QuartoInternalColor2}{ compare  }\textcolor{QuartoInternalColor6}{# isort:skip # noqa}\textcolor{QuartoInternalColor2}{}
\textcolor{QuartoInternalColor2}{}\textcolor{QuartoInternalColor4}{     12}\textcolor{QuartoInternalColor2}{ }\textcolor{QuartoInternalColor5}{from}\textcolor{QuartoInternalColor2}{}\textcolor{QuartoInternalColor7}{ }\textcolor{QuartoInternalColor2}{}\textcolor{QuartoInternalColor2}{ydata\_profiling}\textcolor{QuartoInternalColor2}{}\textcolor{QuartoInternalColor2}{.}\textcolor{QuartoInternalColor2}{}\textcolor{QuartoInternalColor2}{controller}\textcolor{QuartoInternalColor2}{}\textcolor{QuartoInternalColor7}{ }\textcolor{QuartoInternalColor2}{}\textcolor{QuartoInternalColor5}{import}\textcolor{QuartoInternalColor2}{ pandas\_decorator  }\textcolor{QuartoInternalColor6}{# isort:skip # noqa}\textcolor{QuartoInternalColor2}{}
\textcolor{QuartoInternalColor2}{}\textcolor{QuartoInternalColor4}{     13}\textcolor{QuartoInternalColor2}{ }\textcolor{QuartoInternalColor5}{from}\textcolor{QuartoInternalColor2}{}\textcolor{QuartoInternalColor7}{ }\textcolor{QuartoInternalColor2}{}\textcolor{QuartoInternalColor2}{ydata\_profiling}\textcolor{QuartoInternalColor2}{}\textcolor{QuartoInternalColor2}{.}\textcolor{QuartoInternalColor2}{}\textcolor{QuartoInternalColor2}{profile\_report}\textcolor{QuartoInternalColor2}{}\textcolor{QuartoInternalColor7}{ }\textcolor{QuartoInternalColor2}{}\textcolor{QuartoInternalColor5}{import}\textcolor{QuartoInternalColor2}{ ProfileReport  }\textcolor{QuartoInternalColor6}{# isort:skip # noqa}\textcolor{QuartoInternalColor2}{}
\textcolor{QuartoInternalColor2}{}\textcolor{QuartoInternalColor3}{File }\textcolor{QuartoInternalColor2}{}\textcolor{QuartoInternalColor4}{\textasciitilde{}/Desktop/master/2sem/LCPB/projects/.venv/lib/python3.13/site-packages/ydata\_profiling/compare\_reports.py:12}\textcolor{QuartoInternalColor2}{}
\textcolor{QuartoInternalColor2}{}\textcolor{QuartoInternalColor4}{     10}\textcolor{QuartoInternalColor2}{ }\textcolor{QuartoInternalColor5}{from}\textcolor{QuartoInternalColor2}{}\textcolor{QuartoInternalColor7}{ }\textcolor{QuartoInternalColor2}{}\textcolor{QuartoInternalColor2}{ydata\_profiling}\textcolor{QuartoInternalColor2}{}\textcolor{QuartoInternalColor2}{.}\textcolor{QuartoInternalColor2}{}\textcolor{QuartoInternalColor2}{model}\textcolor{QuartoInternalColor2}{}\textcolor{QuartoInternalColor7}{ }\textcolor{QuartoInternalColor2}{}\textcolor{QuartoInternalColor5}{import}\textcolor{QuartoInternalColor2}{ BaseDescription}
\textcolor{QuartoInternalColor2}{}\textcolor{QuartoInternalColor4}{     11}\textcolor{QuartoInternalColor2}{ }\textcolor{QuartoInternalColor5}{from}\textcolor{QuartoInternalColor2}{}\textcolor{QuartoInternalColor7}{ }\textcolor{QuartoInternalColor2}{}\textcolor{QuartoInternalColor2}{ydata\_profiling}\textcolor{QuartoInternalColor2}{}\textcolor{QuartoInternalColor2}{.}\textcolor{QuartoInternalColor2}{}\textcolor{QuartoInternalColor2}{model}\textcolor{QuartoInternalColor2}{}\textcolor{QuartoInternalColor2}{.}\textcolor{QuartoInternalColor2}{}\textcolor{QuartoInternalColor2}{alerts}\textcolor{QuartoInternalColor2}{}\textcolor{QuartoInternalColor7}{ }\textcolor{QuartoInternalColor2}{}\textcolor{QuartoInternalColor5}{import}\textcolor{QuartoInternalColor2}{ Alert}
\textcolor{QuartoInternalColor2}{}\textcolor{QuartoInternalColor4}{---> }\textcolor{QuartoInternalColor2}{}\textcolor{QuartoInternalColor4}{12}\textcolor{QuartoInternalColor2}{ }\textcolor{QuartoInternalColor5}{from}\textcolor{QuartoInternalColor2}{}\textcolor{QuartoInternalColor7}{ }\textcolor{QuartoInternalColor2}{}\textcolor{QuartoInternalColor2}{ydata\_profiling}\textcolor{QuartoInternalColor2}{}\textcolor{QuartoInternalColor2}{.}\textcolor{QuartoInternalColor2}{}\textcolor{QuartoInternalColor2}{profile\_report}\textcolor{QuartoInternalColor2}{}\textcolor{QuartoInternalColor7}{ }\textcolor{QuartoInternalColor2}{}\textcolor{QuartoInternalColor5}{import}\textcolor{QuartoInternalColor2}{ ProfileReport}
\textcolor{QuartoInternalColor2}{}\textcolor{QuartoInternalColor4}{     15}\textcolor{QuartoInternalColor2}{ }\textcolor{QuartoInternalColor5}{def}\textcolor{QuartoInternalColor2}{}\textcolor{QuartoInternalColor7}{ }\textcolor{QuartoInternalColor2}{}\textcolor{QuartoInternalColor8}{\_should\_wrap}\textcolor{QuartoInternalColor2}{(v1: Any, v2: Any) -> }\textcolor{QuartoInternalColor5}{bool}\textcolor{QuartoInternalColor2}{:}
\textcolor{QuartoInternalColor2}{}\textcolor{QuartoInternalColor4}{     16}\textcolor{QuartoInternalColor2}{     }\textcolor{QuartoInternalColor5}{if}\textcolor{QuartoInternalColor2}{ }\textcolor{QuartoInternalColor5}{isinstance}\textcolor{QuartoInternalColor2}{(v1, (}\textcolor{QuartoInternalColor5}{list}\textcolor{QuartoInternalColor2}{, }\textcolor{QuartoInternalColor5}{dict}\textcolor{QuartoInternalColor2}{)):}
\textcolor{QuartoInternalColor2}{}\textcolor{QuartoInternalColor3}{File }\textcolor{QuartoInternalColor2}{}\textcolor{QuartoInternalColor4}{\textasciitilde{}/Desktop/master/2sem/LCPB/projects/.venv/lib/python3.13/site-packages/ydata\_profiling/profile\_report.py:11}\textcolor{QuartoInternalColor2}{}
\textcolor{QuartoInternalColor2}{}\textcolor{QuartoInternalColor4}{      9}\textcolor{QuartoInternalColor2}{ }\textcolor{QuartoInternalColor5}{with}\textcolor{QuartoInternalColor2}{ warnings.catch\_warnings():}
\textcolor{QuartoInternalColor2}{}\textcolor{QuartoInternalColor4}{     10}\textcolor{QuartoInternalColor2}{     warnings.simplefilter(}\textcolor{QuartoInternalColor9}{"}\textcolor{QuartoInternalColor2}{}\textcolor{QuartoInternalColor9}{ignore}\textcolor{QuartoInternalColor2}{}\textcolor{QuartoInternalColor9}{"}\textcolor{QuartoInternalColor2}{)}
\textcolor{QuartoInternalColor2}{}\textcolor{QuartoInternalColor4}{---> }\textcolor{QuartoInternalColor2}{}\textcolor{QuartoInternalColor4}{11}\textcolor{QuartoInternalColor2}{     }\textcolor{QuartoInternalColor5}{import}\textcolor{QuartoInternalColor2}{}\textcolor{QuartoInternalColor7}{ }\textcolor{QuartoInternalColor2}{}\textcolor{QuartoInternalColor2}{pkg\_resources}\textcolor{QuartoInternalColor2}{}
\textcolor{QuartoInternalColor2}{}\textcolor{QuartoInternalColor4}{     13}\textcolor{QuartoInternalColor2}{ }\textcolor{QuartoInternalColor5}{if}\textcolor{QuartoInternalColor2}{ }\textcolor{QuartoInternalColor10}{not}\textcolor{QuartoInternalColor2}{ is\_pyspark\_installed():}
\textcolor{QuartoInternalColor2}{}\textcolor{QuartoInternalColor4}{     14}\textcolor{QuartoInternalColor2}{     }\textcolor{QuartoInternalColor5}{from}\textcolor{QuartoInternalColor2}{}\textcolor{QuartoInternalColor7}{ }\textcolor{QuartoInternalColor2}{}\textcolor{QuartoInternalColor2}{typing}\textcolor{QuartoInternalColor2}{}\textcolor{QuartoInternalColor7}{ }\textcolor{QuartoInternalColor2}{}\textcolor{QuartoInternalColor5}{import}\textcolor{QuartoInternalColor2}{ TypeVar}
\textcolor{QuartoInternalColor2}{}\textcolor{QuartoInternalColor1}{ModuleNotFoundError}\textcolor{QuartoInternalColor2}{: No module named 'pkg\_resources'}
\end{Highlighting}

}

\end{codelisting}%

\subsection{Preprocessing}\label{preprocessing}

\begin{itemize}
\tightlist
\item
  Event cleaning
\item
  Time ordering
\item
  Normalization of features
\end{itemize}

\begin{center}\rule{0.5\linewidth}{0.5pt}\end{center}

\section{Feature Engineering}\label{feature-engineering}

\subsection{Event Pair Construction}\label{event-pair-construction}

Candidate signal events are constructed by pairing events based on
temporal proximity.

For each pair, the following features are computed:

\begin{itemize}
\tightlist
\item
  Time difference \(\Delta t\)
\item
  Spatial distance \(\Delta r\)
\item
  Energy of prompt and delayed signals
\item
  Energy correlation
\end{itemize}

\subsection{Feature Selection}\label{feature-selection}

Relevant features are selected based on their ability to discriminate
between signal and background.

\begin{center}\rule{0.5\linewidth}{0.5pt}\end{center}

\section{Machine Learning Model}\label{machine-learning-model}

\subsection{Problem Formulation}\label{problem-formulation}

We define a binary classification task:

\begin{itemize}
\tightlist
\item
  \(y = 1\): signal
\item
  \(y = 0\): background
\end{itemize}

\subsection{Model Architecture}\label{model-architecture}

Possible models:

\begin{itemize}
\tightlist
\item
  Feedforward Neural Network
\item
  Gradient Boosted Trees
\end{itemize}

\subsection{Training Procedure}\label{training-procedure}

\begin{itemize}
\tightlist
\item
  Train on balanced simulated dataset
\item
  Use train/validation split
\item
  Optimize classification performance
\end{itemize}

\begin{center}\rule{0.5\linewidth}{0.5pt}\end{center}

\section{Model Calibration}\label{model-calibration}

\subsection{Motivation}\label{motivation}

The training dataset is artificially balanced, while real data is highly
imbalanced, leading to miscalibrated outputs .

\subsection{Calibration Methods}\label{calibration-methods}

\begin{itemize}
\tightlist
\item
  Platt scaling
\item
  Isotonic regression
\item
  Likelihood reweighting
\end{itemize}

\begin{center}\rule{0.5\linewidth}{0.5pt}\end{center}

\section{Event Selection Performance}\label{event-selection-performance}

\subsection{Metrics}\label{metrics}

\begin{itemize}
\tightlist
\item
  Efficiency (signal acceptance)
\item
  Purity (signal fraction)
\item
  ROC curve and AUC
\end{itemize}

\subsection{Results}\label{results}

\begin{itemize}
\tightlist
\item
  Performance comparison across models
\item
  Optimal selection threshold
\end{itemize}

\begin{center}\rule{0.5\linewidth}{0.5pt}\end{center}

\section{Energy Spectrum
Reconstruction}\label{energy-spectrum-reconstruction}

\subsection{Event Selection}\label{event-selection}

Apply the trained and calibrated model to select neutrino candidates.

\subsection{Spectrum Construction}\label{spectrum-construction}

Reconstruct the neutrino energy spectrum from selected events.

\begin{center}\rule{0.5\linewidth}{0.5pt}\end{center}

\section{Oscillation Analysis}\label{oscillation-analysis}

\subsection{Near vs Far Detector
Comparison}\label{near-vs-far-detector-comparison}

Oscillation effects are extracted by comparing spectra between
detectors, reducing systematic uncertainties .

\subsection{Parameter Extraction}\label{parameter-extraction}

Fit oscillation parameters from the observed spectrum distortion.

\begin{center}\rule{0.5\linewidth}{0.5pt}\end{center}

\section{Discussion}\label{discussion}

\subsection{Limitations}\label{limitations}

\begin{itemize}
\tightlist
\item
  Background dominance
\item
  Model dependence on simulation
\item
  Calibration uncertainties
\end{itemize}

\subsection{Improvements}\label{improvements}

\begin{itemize}
\tightlist
\item
  More robust feature engineering
\item
  Domain adaptation techniques
\item
  Better uncertainty estimation
\end{itemize}

\begin{center}\rule{0.5\linewidth}{0.5pt}\end{center}

\section{Conclusion}\label{conclusion}

We developed a machine learning pipeline for neutrino event selection in
a high-background environment. The method successfully identifies signal
events, reconstructs the energy spectrum, and enables the extraction of
oscillation parameters.

\begin{center}\rule{0.5\linewidth}{0.5pt}\end{center}


\renewcommand\refname{References}
\bibliography{references.bib}



\end{document}
